\documentclass[11pt]{article}

\usepackage{amsmath,amssymb}
\usepackage{enumitem}
\usepackage{geometry}
\geometry{margin=1in}

\title{Sample Oral Exam Questions and Answers\\
Immersed Boundary Method: Integral Transforms and Regularized Delta Functions}
\author{}
\date{}

\begin{document}
\maketitle

Throughout, $\mathbf{q}\in U$ is a fixed material label and
$\mathbf{X}(\mathbf{q},t)\in\Omega$ is its current position. This is the Peskin
convention used in Module~1; Module~4 (IBFE) switches to continuum-mechanics
notation $\boldsymbol{\chi}(\mathbf{X},t)$ with $\mathbf{X}$ the reference point.
Lagrangian force densities are understood with respect to the explicitly stated
reference measure. These conventions determine both units and Jacobian factors.

\section*{Eulerian--Lagrangian Framework}

\begin{enumerate}[label=\arabic*.]

\item \textbf{What are the Eulerian and Lagrangian variables in the immersed boundary method?}

\textit{Sample answer:}  
The fluid is described in Eulerian coordinates by velocity
$\mathbf{u}(\mathbf{x},t)$ and pressure $p(\mathbf{x},t)$ on the fixed spatial
domain $\Omega$. The structure is described in material coordinates by the map
$\mathbf{X}(\mathbf{q},t)$, which gives the current physical position of
the material label $\mathbf{q}$ in a reference domain $U$.

\item \textbf{Why are integral transforms needed to couple Eulerian and Lagrangian variables?}

\textit{Sample answer:}  
Eulerian and Lagrangian fields live on different domains and usually different
meshes. A delta distribution defines the push-forward of a Lagrangian force to
Eulerian space and the pullback (evaluation along the motion) of Eulerian
velocity to the structure. At the discrete level, regularized kernels and
quadrature approximate this dual pairing.

\item \textbf{Write the Eulerian--Lagrangian coupling equations used in IB.}

\textit{Sample answer:}  
Force spreading:
\[
\mathbf{f}(\mathbf{x},t) = \int_U \mathbf{F}(\mathbf{q},t)\,
\delta\!\left(\mathbf{x}-\mathbf{X}(\mathbf{q},t)\right)\,d\mathbf{q}.
\]
Velocity interpolation:
\[
\frac{\partial \mathbf{X}}{\partial t}(\mathbf{q},t)
= \int_\Omega \mathbf{u}(\mathbf{x},t)\,
\delta\!\left(\mathbf{x}-\mathbf{X}(\mathbf{q},t)\right)\,d\mathbf{x}.
\]
The second equation is the no-slip kinematic relation for the classical
massless IB structure; a structure with independent inertia requires its own
momentum equation and a corresponding coupling law.

\item \textbf{What is the physical meaning of the mapping $\mathbf{X} : U \times [0,T] \to \Omega$?}

\textit{Sample answer:}
The map $\mathbf{X}$ takes a fixed material label in a chosen reference
configuration to its current spatial position. It encodes deformation and
rigid motion. The reference need not be stress-free or literally undeformed;
that is an additional modeling choice. For a volumetric solid, the deformation
gradient and its determinant,
\[
\mathbf{F}_{\rm def}(\mathbf{q},t)=\frac{\partial\mathbf{X}}{\partial\mathbf{q}},
\qquad
J=\det\mathbf{F}_{\rm def},
\]
describe local deformation. The subscript matters: $\mathbf{F}_{\rm def}$ is the
deformation gradient, whereas the unsubscripted $\mathbf{F}$ denotes the
Lagrangian force density elsewhere in these questions.

\item \textbf{Why are two Lagrangian notations commonly used in IB: $\mathbf{X}(\mathbf{q},t)$ and $\boldsymbol{\chi}(\mathbf{X},t)$?}

\textit{Sample answer:}
In classical IB (Peskin) notation, used in Module~1, $\mathbf{q}$ is a fixed
material label and $\mathbf{X}(\mathbf{q},t)$ denotes the current position. In
continuum-mechanics notation, used from Module~4 on, $\mathbf{X}$ is instead
reserved for the fixed reference coordinate and $\boldsymbol{\chi}(\mathbf{X},t)$
is the current-position map. Both are valid,
but one must not let the same symbol mean both a material label and a current
position in one derivation.

\end{enumerate}

\section*{Reference Configuration and Force Definitions}

\begin{enumerate}[label=\arabic*., resume]

\item \textbf{What is the Lagrangian reference domain $U$ and why is it fixed in time?}

\textit{Sample answer:}  
$U$ is a chosen material reference configuration. It is fixed because its
coordinates label material points independently of their subsequent motion.
It may be undeformed and stress-free, but neither property follows merely from
using a reference description.

\item \textbf{Explain the distinction between $U$, $\mathbf{X}(U,t)$, and $\Omega$.}

\textit{Sample answer:}
$U$ is the material reference domain,
$\mathbf{X}(U,t)\subseteq\Omega$ is the current structural configuration,
and $\Omega$ is the Eulerian spatial domain. Integrals over these sets use
different measures. A Jacobian is required when changing between reference and
current measures, but not when an integral is deliberately kept over $U$.

\item \textbf{What is the first Piola--Kirchhoff stress and why is it used in IBFE?}

\textit{Sample answer:}  
The first Piola--Kirchhoff stress $\mathbf{P}_e$ maps a reference-area normal to
a force measured in the current configuration per unit reference area. For a
hyperelastic energy density $W(\mathbf{F}_{\rm def})$ per reference volume,
$\mathbf{P}_e=\partial W/\partial\mathbf{F}_{\rm def}$. It is generally nonsymmetric and is
work-conjugate to the deformation gradient, which makes it natural in IBFE's
reference-domain weak form.

\item \textbf{How is the Lagrangian elastic force density defined?}

\textit{Sample answer:}  
In the interior of a smooth volumetric solid, one common sign convention gives
\[
\mathbf{F}_e(\mathbf{q},t) = \nabla_{\!q} \cdot \mathbf{P}_e(\mathbf{q},t),
\qquad
(\nabla_{\!q}\cdot\mathbf{P}_e)_i=\frac{\partial (P_e)_{iJ}}{\partial q_J},
\]
which is force per unit reference volume. Note that the material divergence
$\nabla_{\!q}\cdot$ contracts on the \emph{reference} (second) index of
$\mathbf{P}_e$; the first index is a current-configuration index, so no
contraction over it would be meaningful. Boundary tractions are also required:
if elastic force is minus the variation of stored energy, the associated
reference-boundary force is $-\mathbf{P}_e\mathbf{N}$ under this convention.
IBFE often avoids forming either strong term and defines the force by
$\int_U\mathbf{F}_e\cdot\mathbf{V}\,d\mathbf{q}
=-\int_U\mathbf{P}_e:\nabla_{\!q}\mathbf{V}\,d\mathbf{q}$, with boundary conditions
handled consistently.

\end{enumerate}

\section*{Force Spreading and the Absence of a Jacobian}

\begin{enumerate}[label=\arabic*., resume]

\item \textbf{Write the classical immersed boundary force-spreading formula.}

\textit{Sample answer:}
\[
\mathbf{f}(\mathbf{x},t) = \int_U \mathbf{F}(\mathbf{q},t)\,
\delta\!\left(\mathbf{x}-\mathbf{X}(\mathbf{q},t)\right)\,d\mathbf{q}.
\]

\item \textbf{Why does no Jacobian appear in this integral?}

\textit{Sample answer:}
Because $\mathbf{F}$ is defined per unit reference measure and the integration
variable remains $\mathbf{q}\in U$. No change of variables has been made. The
geometry enters through $\mathbf{X}(\mathbf{q},t)$ inside the delta distribution, whose
units also ensure that the resulting $\mathbf{f}$ has the correct Eulerian
force-density units.

\item \textbf{Why is force spreading not a change of variables?}

\textit{Sample answer:}  
A change of variables rewrites the same integral using a new measure. Spreading
instead defines a push-forward distribution. For a smooth test field
$\boldsymbol{\varphi}$,
\[
\int_\Omega\mathbf{f}\cdot\boldsymbol{\varphi}\,d\mathbf{x}
=\int_U\mathbf{F}\cdot
\boldsymbol{\varphi}(\mathbf{X}(\mathbf{q},t))\,d\mathbf{q}.
\]
This defining identity contains no coordinate substitution.

\item \textbf{Under what conditions would a Jacobian appear?}

\textit{Sample answer:}  
A Jacobian appears when a current-configuration density is rewritten using a
reference measure (or conversely). For a volumetric solid,
$d\mathbf{x}=J\,d\mathbf{q}$; surface and curve measures use the corresponding
surface or line Jacobians. The factor therefore depends on how the force density
was defined, not merely on whether the object is a solid, surface, or fiber.

\end{enumerate}

\section*{Physical Interpretation and Conservation}

\begin{enumerate}[label=\arabic*., resume]

\item \textbf{Interpret force spreading by integrating over a fluid region $\Gamma \subset \Omega$.}

\textit{Sample answer:}  
For the exact delta distribution,
\[
\int_\Gamma\mathbf{f}(\mathbf{x},t)\,d\mathbf{x}
=\int_U\mathbf{F}(\mathbf{q},t)\,
\mathbf{1}_{\Gamma}(\mathbf{X}(\mathbf{q},t))\,d\mathbf{q},
\]
up to the usual convention for points on $\partial\Gamma$. With a regularized
kernel, points within one kernel radius of $\partial\Gamma$ contribute only the
fraction of the kernel support lying in $\Gamma$.

\item \textbf{Why does IB conserve total force?}

\textit{Sample answer:}  
In the continuum, $\int_\Omega\delta(\mathbf{x}-\mathbf{X}(\mathbf{q},t))\,d\mathbf{x}=1$
when the support is contained in the domain, so total Eulerian and Lagrangian
forces agree. Discretely, the analogous result requires the kernel's zeroth
moment (partition-of-unity) condition, consistent quadrature weights, and
appropriate treatment near physical or periodic boundaries.

\item \textbf{Why is the Eulerian force density a distribution?}

\textit{Sample answer:}  
For a codimension-one surface or codimension-two fiber, the force is supported
on a set of zero Eulerian volume, so its Eulerian representation is a singular
distribution. A force spread from a volumetric immersed solid need not be
singular in its occupied volume, although jumps at its boundary can still be
represented distributionally. Numerical regularization replaces these singular
objects by grid-scale functions.

\item \textbf{Under what conditions is the velocity continuous across the structure while the pressure jumps?}

\textit{Sample answer:}
Velocity is continuous whenever the classical no-slip kinematic coupling holds
and the interface carries no singular mass or slip source; the pressure jumps
only when the interfacial force has a nonzero normal component, so neither
property is automatic. A singular interfacial force produces
a traction jump; its normal component is often associated with a pressure jump,
while tangential components can produce viscous-stress jumps. Pressure need not
jump for every force distribution, and a regularized discrete method smooths
all such jumps over the kernel width. Viscosity alone is not the kinematic
constraint.

\end{enumerate}

\section*{Regularized Delta Functions}

\begin{enumerate}[label=\arabic*., resume]

\item \textbf{Why must the Dirac delta be regularized numerically?}

\textit{Sample answer:}  
The Dirac delta has no finite point values and cannot be sampled as an ordinary
grid function. A discrete regularized delta replaces it by a compact stencil
whose width scales with the mesh. Conservation and interpolation properties are
obtained only if the kernel satisfies the required discrete moments and is used
with consistent quadrature and boundary treatment.

\item \textbf{Why must the same kernel be used for spreading and interpolation?}

\textit{Sample answer:}  
Using the same kernel is part of making interpolation the adjoint of spreading.
With matching Eulerian and Lagrangian quadrature weights, this gives the power
identity
$\langle\mathbf{u},\mathcal{S}\mathbf{F}\rangle_E
=\langle\mathcal{J}\mathbf{u},\mathbf{F}\rangle_L$.
Using the same formula without consistent weights or inner products is not, by
itself, sufficient.

\item \textbf{What does it mean that force spreading is a local averaging operation?}

\textit{Sample answer:}  
Each Lagrangian quadrature force is distributed over the nearby grid points in
the compact kernel support. Conversely, interpolation forms a kernel-weighted
local average of grid velocity. This locality controls stencil cost and
regularization width; it should not be confused with averaging in time or over
an ensemble.

\end{enumerate}

\section*{Design Criteria for Regularized Delta Functions}

\begin{enumerate}[label=\arabic*., resume]

\item \textbf{State Peskin’s zeroth and first moment conditions and their physical meaning.}

\textit{Sample answer:}  
In one dimension, for every real shift $r$, representative discrete conditions
are
\[
\sum_{j\in\mathbb{Z}}\phi(r-j)=1,
\qquad
\sum_{j\in\mathbb{Z}}(r-j)\phi(r-j)=0.
\]
Tensor-product versions apply in several dimensions. The zeroth moment gives
partition of unity and hence total-force conservation. The first moment
reproduces the force centroid; together with the zeroth moment it preserves the
torque of spread point forces on an unbounded or suitably treated grid.

\item \textbf{Why are symmetry and translational invariance important?}

\textit{Sample answer:}  
Kernel symmetry avoids a preferred direction and supports centered moment
conditions. Exact continuous translational invariance is impossible for a
compact stencil on a fixed grid; discrete moment, even--odd, and sum-of-squares
conditions are used to reduce variation as a marker moves relative to grid
lines. The \emph{even--odd condition} is the requirement that, for every real
shift $r$,
\[
\sum_{j\ \mathrm{even}}\phi(r-j)
=\sum_{j\ \mathrm{odd}}\phi(r-j)=\tfrac12 ;
\]
its purpose is to keep the even- and odd-indexed grid points from decoupling in
the spread and interpolate operators, which would otherwise admit checkerboard
(sawtooth) modes invisible to the interaction. It is distinct from the
sum-of-squares condition $\sum_j\phi(r-j)^2=C$, which instead controls the
approximate translation invariance of the interaction strength. Peskin's
four-point kernel satisfies the even--odd condition; the three-point
Roma--Peskin--Berger kernel does not need it, because a three-cell stencil
always couples adjacent grid points. State which invariance property is meant
rather than claiming exact grid independence.

\item \textbf{What is the tradeoff between smoothness and cost?}

\textit{Sample answer:}  
Additional smoothness can reduce force and velocity jitter as markers cross the
grid, but usually requires wider support. Wider stencils increase work and
communication and thicken the numerical interface; they do not automatically
increase the convergence order of the full IB scheme.

\end{enumerate}

\section*{Common Kernels and Practice}

\begin{enumerate}[label=\arabic*., resume]

\item \textbf{Why is Peskin’s 4-point kernel commonly used?}

\textit{Sample answer:}  
Peskin's classical four-point tensor-product kernel has compact four-cell
support in each coordinate and satisfies a useful collection of discrete
moment, even--odd, and sum-of-squares conditions. It is a well-tested compromise
between grid invariance, smoothness, and stencil cost, but it is not optimal for
every accuracy, leakage, or conservation objective.

\item \textbf{What happens as $h \to 0$?}

\textit{Sample answer:}  
A properly scaled kernel family can converge to the Dirac delta in the
distributional sense: its support shrinks and its moments approach those of a
point mass. Convergence of the \emph{full} discrete IB method additionally
requires consistent spatial operators, Lagrangian quadrature and refinement,
stable time integration, solver convergence, and appropriate treatment of the
regularized interface. Kernel convergence alone is not a convergence proof for
the coupled algorithm.

\end{enumerate}

\end{document}
